

%: Style file for Latex
% Most style definitions are in the external file PhDthesisPSnPDF.
% In this template package, it can be found in ./Latex/Classes/
\documentclass[twoside,12pt]{Latex/Classes/PhDthesisPSnPDF}


%: Macro file for Latex
% Macros help you summarise frequently repeated Latex commands.
% Here, they are placed in an external file /Latex/Macros/MacroFile1.tex
% An macro that you may use frequently is the figuremacro (see introduction.tex)
%---------------------------------------------------------------
% Macros
% version 3 by Igor Ruiz-Agundez 2011
% version 2 by Jakob Suckale 2007
% version 1 by Harish Bhanderi 2002
%---------------------------------------------------------------

% This file contains macros that can be called up from connected TeX files
% It helps to summarise repeated code, e.g. figure insertion (see below).


%---------------------------------------------------------------
% MY COMMANDS
%---------------------------------------------------------------
%MY COMMANDS
\newcommand{\sub}[1]{\mbox{\scriptsize{#1}}}
\newcommand{\der}[2]{ \frac{ \partial #1 }{\partial #2} }
\newcommand{\dtot}[2]{ \frac{ d #1 }{d #2} }
\newcommand{\pr}[1]{ \left( #1 \right) }
\newcommand{\cor}[1]{ \left[ #1 \right] }
\newcommand{\lla}[1]{ \left\{ #1 \right\} }
\newcommand{\eq}[2]{\begin{equation} \label{#1} #2 \end{equation}}
\newcommand{\bds}[1]{\boldsymbol{ #1 }}
\newcommand{\oiint}{\displaystyle\bigcirc\!\!\!\!\!\!\!\!\int\!\!\!\!\!\int}
\newcommand{\mathsize}[2]{\mbox{\fontsize{#1}{#1}\selectfont $#2$}}
\newcommand{\sinc}[1]{\mbox{sinc}#1}

\newcommand{\cita}[1]{\textsuperscript{\tiny\cite{#1}}}
\newcommand{\ket}{\rangle}
\newcommand{\bra}{\langle}

\renewcommand{\listtablename}{Índice de tablas}
\renewcommand{\tablename}{Tabla}

\usepackage{color, colortbl}
%Celda de colores en tablas
\newcommand{\cellc}[1]{\multicolumn{1}{|>{\columncolor[rgb]{0.8, 0.8, 0.8}}c|}{#1}}

%---------------------------------------------------------------
% Figures
%---------------------------------------------------------------


% Makes the \InsertFig macro compatible both with one or two columns
\makeatletter
\newlength \figwidth
\if@twocolumn
  \setlength \figwidth {\columnwidth}
\else
  \setlength \figwidth {\textwidth}
\fi
\makeatother

% \InsertFig allows inserting figures
% Parameters
% 1 --> Filename
% 2 --> Label for referencing
% 3 --> Title describing the figure (caption)
% 4 --> Description of the figure
% 5 --> Figure width, range [0,1]. If parameter is left blank the figure size is not change
% 6 --> Any other option for \includegraphics
% Usage:
% \InsertFig{}{}{}{}{}{}
%
\newcommand{\InsertFig}[6]{%
	\ifthenelse{\isempty{#5}}%
	{% if #1 is empty
		\begin{figure}[htbp!]
		\centering
		\includegraphics[#6]{#1}%
		\caption{#3}{\textbf{#4}}
		\label{#2}
		\end{figure}    
	}
	{% if #1 is not empty
		\begin{figure}[htbp!]
		\centering
		\includegraphics[width=#5\figwidth,#6]{#1}%
		\caption{#3}{\textbf{#4}}
		\label{#2}
		\end{figure}
	}
}

%% Simple version of \InsertFig
%\newcommand{\InsertFig}[5]{
%  \begin{figure}[htbp]
%   	\centering
%    \includegraphics[width=#4\textwidth,#5]{#1}%
%    \caption{#3}
%    \label{#2}
%  \end{figure}
%}



% insert a centered figure with caption
% parameters 1:filename, 2:label, 3:title, 
\newcommand{\figuremacro}[3]{
	\begin{figure}[htbp]
		\centering
		\includegraphics[width=1\textwidth]{#1}
		\caption[#3]{\textbf{#3}}
		\label{#2}
	\end{figure}
}


% insert a centered figure with caption and description
% parameters 1:filename, 2:label, 3:title, 4:description
\newcommand{\figuremacroD}[4]{
	\begin{figure}[htbp]
		\centering
		\includegraphics[width=1\textwidth]{#1}
		\caption[#3]{\textbf{#3} - #4}
		\label{#2}
	\end{figure}
}

% insert a centered figure with caption and description AND WIDTH
% parameters 1:filename, 2:label, 3:title, 4:description, 5: textwidth
% textwidth 1 means as text, 0.5 means half the width of the text
\newcommand{\figuremacroDW}[5]{
	\begin{figure}[htbp]
		\centering
		\includegraphics[width=#5\textwidth]{#1}
		\caption[#3]{\textbf{#3} - #4}
		\label{#2}
	\end{figure}
}

% inserts a figure with wrapped around text; only suitable for NARROW figs
% o is for outside on a double paged document; others: l, r, i(inside)
% text and figure will each be half of the document width
% note: long captions often crash with adjacent content; take care
% in general: above 2 macro produce more reliable layout
\newcommand{\figuremacroN}[3]{
	\begin{wrapfigure}{o}{0.5\textwidth}
		\centering
		\includegraphics[width=0.48\textwidth]{#1}
		\caption[#2]{{\small\textbf{#2} - #3}}
		\label{#1}
	\end{wrapfigure}
}




% Estas definiciones son para el comando \InsertFigBox
\newlength{\anchoFigura}
\newlength{\anchoFloat}
\addtolength{\fboxsep}{2\fboxsep}
%\renewcommand{\capfont}{\normalfont\normalcolor\sffamily\small}
%\renewcommand{\caplabelfont}{\normalfont\normalcolor\sffamily\bfseries\small}

% El comando \InsertFigBox nos permite insertar figuras en un marco
% Los parametros son:
% 1 --> Fichero de la imagen
% 2 --> Etiqueta (label) para referencias
% 3 --> Texto a pie de imagen
% 4 -> Porcentaje del ancho de página que ocupará la figura (de 0 a 1)
% 5 --> Opciones que queramos pasarle al \includegraphics
\newcommand{\InsertFigBox}[5]{%
  \setlength{\anchoFloat}{#4\textwidth}%
  \addtolength{\anchoFloat}{-4\fboxsep}%
  \setlength{\anchoFigura}{\anchoFloat}%
  \begin{figure}%
    \begin{center}%
      \Ovalbox{%
        \begin{minipage}{\anchoFloat}%
          \begin{center}%
            \includegraphics[width=\anchoFigura,#5]{#1}%
            \caption{#3}%
            \label{#2}%
          \end{center}%
        \end{minipage}
      }%
    \end{center}%
  \end{figure}%
}



%---------------------------------------------------------------
% Misc
%---------------------------------------------------------------

% predefined commands by Harish
\newcommand{\PdfPsText}[2]{
  \ifpdf
     #1
  \else
     #2
  \fi
}


%---------------------------------------------------------------
% Locales
%---------------------------------------------------------------


%%
%% Para quitar traducciones raras (Cuadros)
%% A de usarse cada vez que se seleccione el idioma
%%
\newcommand{\MejorarTraducciones}{%
       \renewcommand{\listtablename}{Índice de tablas}
       \renewcommand{\tablename}{Tabla}
       \renewcommand{\lstlistingname}{Lista}
}%



%---------------------------------------------------------------
% Source code
%---------------------------------------------------------------


%%
%% Para escribir extractos de codigo
%%
%% Las tabulaciones se substituyen por dos espacios
%\fvset{tabsize=2}
%% Creamos un nuevo environment de fancyvrb para los ejemplos enmarcados
%\DefineVerbatimEnvironment{VerbEj}{BVerbatim}{fontsize=\small,samepage=true,commandchars=\\\{\}}
%% Colo de fondo
%\definecolor{grisfondo}{gray}{0.9}
%% Environment para extractos de codigo
%\newenvironment{codigo}%
%{\VerbatimEnvironment\begin{Sbox}\begin{VerbEj}}%
%{\end{VerbEj}\end{Sbox}\setlength{\fboxsep}{8pt}\begin{center}\fcolorbox{black}{grisfondo}{\TheSbox}\end{center}}
%
%% Otro formato más bonito para código fuente
%\newcommand{\codigofuente}[3]{%
%  \lstinputlisting[language=#1,caption={#2}]{#3}%
%}






%: ----------------------------------------------------------------------
%:                  TITLE PAGE: name, degree,..
% ----------------------------------------------------------------------
% below is to generate the title page with crest and author name

% if output to PDF then put the following in PDF header
\ifpdf  
    \pdfinfo { /Title  (PhD)
               /Creator (TeX)
               /Producer (pdfTeX)
               /Author (Name surname)
               /CreationDate (D:YYYYMMDDhhmmss)  %format D:YYYYMMDDhhmmss
               /ModDate (D:YYYYMMDDhhmm)
               /Subject (xyz)
               /Keywords (keyword1, keyword2, keyword3) }
    \pdfcatalog { /PageMode (/UseOutlines)
                  /OpenAction (fitbh)  }
\fi


% Title of the dissertation
\title{The place of the Milky Way and Andromeda in the cosmic web}


% ----------------------------------------------------------------------
% This section below defines front covert (external and internal)
% Shield logo
\crest{\includegraphics[width=3cm]{UdeA_Shield}}
% Full logo
%\crest{\includegraphics[width=6cm]{UDeusto}}
\university{Universidad de Antioquia \\ Facultad de Ciencias Exactas y Naturales \\ Instituto de Física                                  }
\degree{}
\author{\textbf{Sebastian Bustamante Jaramillo}} 
\collegeordept{Facultad de Ciencias Exactas y Naturales \\ Instituto de Física}
\textadvisor{Advisor: \\}
\advisor{\textbf{Prof. Jaime E. Forero-Romero}}
\textsignaturecandidate{Student}
\textsignatureadvisor{Advisor}
\cityofbirth{Medellín}
%\degreedate{\monthname \ \the\year}
\degreedate{January \the\year}
% ----------------------------------------------------------------------
% turn of those nasty overfull and underfull hboxes
\hbadness=10000
\hfuzz=50pt


%: --------------------------------------------------------------
%:                  FRONT MATTER: dedications, abstract,..
% --------------------------------------------------------------

\begin{document}

\selectlanguage{british}

% sets line spacing
\renewcommand\baselinestretch{1.2}
\baselineskip=18pt plus1pt

% Watermark
%\watermark{DRAFT	DRAFT	DRAFT	DRAFT	DRAFT	DRAFT	DRAFT	DRAFT	DRAFT}


%: ----------------------- generate cover page ------------------------

\maketitle  % command to print the title page with above variables

% Title back
\include{0_frontmatter/title_back}





%: ----------------------- abstract ------------------------

% Your institution may have specific regulations if you need an abstract and where it is to be placed in the document. The default here is just after title.


% The original template provides and abstractseparate environment, if your institution requires them to be separate. I think it's easier to print the abstract from the complete thesis by restricting printing to the relevant page.
% \begin{abstractseparate}
%   \input{Abstract/abstract}
% \end{abstractseparate}


%: ----------------------- tie in front matter ------------------------

% The frontmatter text starts here
\frontmatter

\include{0_frontmatter/dedication}

\include{0_frontmatter/abstract}

\include{0_frontmatter/acknowledgement}

% As abstract contains various languages we set the main language again
\selectlanguage{british}


%: ----------------------- contents ------------------------

\setcounter{secnumdepth}{5} % organisational level that receives a numbers
\setcounter{tocdepth}{5}    % print table of contents for level 3


%%You can also add extra lines to the ToC or to force extra unnumbered section headings to be included. For example, if you wanted to add an entry called Preface, and you didn't want the Preface to be numbered, you'd use these commands:
%\ subsection*{Preface}
%\addcontentsline{toc}{subsection}{Preface} 

\tableofcontents            % print the table of contents
% levels are: 0 - chapter, 1 - section, 2 - subsection, 3 - subsection

%: ----------------------- list of figures/tables ------------------------

\listoffigures	% print list of figures
\listoftables  % print list of tables


%: ----------------------- glossary ------------------------

% Tie in external source file for definitions: /0_frontmatter/glossary.tex
% Glossary entries can also be defined in the main text. See glossary.tex
\include{0_frontmatter/glossary}

%\begin{multicols}{2} % \begin{multicols}{#columns}[header text][space]
%\begin{footnotesize} % scriptsize(7) < footnotesize(8) < small (9) < normal (10)


\label{sec:glossary} % target name for links to glossary

%\end{footnotesize}
%\end{multicols}




%: --------------------------------------------------------------
%:                  MAIN DOCUMENT SECTION
% --------------------------------------------------------------

% the main text starts here with the introduction, 1st chapter,...
\mainmatter

%\renewcommand{\chaptername}{} % uncomment to print only "1" not "Chapter 1"
\pagestyle{fancy}

%: ----------------------- subdocuments ------------------------

% Parts of the thesis are included below. Rename the files as required.
% But take care that the paths match. You can also change the order of appearance by moving the include commands.

%------------------------- introduction ------------------------
\include{1_introduction/0_index}

%---------------------- theoretical frame -----------------------
\include{2_theoretical_frame/0_index}	

%--------------------- Nbody simulations ------------------------
\include{3_nbody_simulations/0_index}
			
%--------------------------- results ----------------------------
\include{4_results/0_index}

% --------------------------------------------------------------
%:                  BACK MATTER: appendices, refs,..
% --------------------------------------------------------------

% the back matter: appendix and references close the thesis
\backmatter


%: ----------------------- appendix ------------------------

%\appendix


%: ----------------------- bibliography ------------------------



% The section below defines how references are listed and formatted
% The default below is 2 columns, small font, complete author names.
% Entries are also linked back to the page number in the text and to external URL if provided in the BibTex file.


% Original version:

% PhDbiblio-url2 = names small caps, title bold & hyperlinked, link to page 
%\begin{multicols}{2} % \begin{multicols}{ # columns}[ header text][ space]
%\begin{tiny} % tiny(5) < scriptsize(7) < footnotesize(8) < small (9)
%
%\bibliographystyle{Latex/Classes/PhDbiblio-url2} % Title is link if provided
%\renewcommand{\bibname}{References} % changes the header; default: Bibliography
%
%\bibliography{9_backmatter/references} % adjust this to fit your BibTex file
%
%\end{tiny}
%\end{multicols}



% Show all bibliography entries
%\nocite*



% If we want bibliography backreference, use unsrt first and the desidered one after

%\bibliographystyle{unsrt} % Defines the bibliography style

%\bibliographystyle{alpha} % Defines the bibliography style

%\bibliographystyle{apa-good} % Defines the bibliography style
%\bibliographystyle{natbib} % Defines the bibliography style

%\bibliographystyle{plainurl}

%\renewcommand{\bibname}{References} % changes the header; default: Bibliography

%To include the references/works cited/bibliography in your Table of Contents, right before the bibliography command, use the command
%\addcontentsline{toc}{section}{References}
%
\bibliographystyle{abbrv}
\bibliography{9_backmatter/references} % adjust this to fit your BibTex file


%\printnomenclature[1.5cm] % [] = distance between entry and description

\printnomenclature % [] = distance between entry and description

% --------------------------------------------------------------
% Various bibliography styles exit. Replace above style as desired.

% in-text refs: (1) (1; 2)
% ref list: alphabetical; author(s) in small caps; initials last name; page(s)
%\bibliographystyle{Latex/Classes/PhDbiblio-case} % title forced lower case
%\bibliographystyle{Latex/Classes/PhDbiblio-bold} % title as in bibtex but bold
%\bibliographystyle{Latex/Classes/PhDbiblio-url} % bold + www link if provided

%\bibliographystyle{Latex/Classes/jmb} % calls style file jmb.bst
% in-text refs: author (year) without brackets
% ref list: alphabetical; author(s) in normal font; last name, initials; page(s)

%\bibliographystyle{plainnat} % calls style file plainnat.bst
% in-text refs: author (year) without brackets
% (this works with package natbib)


% --------------------------------------------------------------


%: Declaration of originality
\include{9_backmatter/declaration}





\end{document}
